\section{Introduction}

Code obfuscation is widely used to make software difficult to inspect
while preserving its intended functionality. In benign settings,
developers employ obfuscation to protect intellectual property,
implementation details, and proprietary business logic. In adversarial
settings, however, malware authors routinely exploit the same
techniques to conceal malicious behavior and frustrate program
analysis and reverse engineering. Common transformations include
identifier renaming, control-flow restructuring, dead-code insertion,
expression rewriting, call indirection, etc~\cite{collberg1997taxonomy}. Although such
transformations are intended to preserve program semantics, they can
substantially obscure how those semantics are expressed in the source
code, thereby increasing the difficulty of understanding the
behavior of malware code~\cite{collberg1997taxonomy,collberg2002manufacturing,
Obfuscation-The-Hidden-Malware,
An-Observational-Investigation-of-Reverse-Engineers-Process-and-Mental-Models}.

This difficulty is well established for human analysts. Nguyen
{\em et al.}~\cite{nguyen2026effectcodeobfuscationhuman}, for example,
showed that code obfuscation significantly reduces humans' accuracy in
reasoning about program outputs. Such findings are particularly
important for malware analysis, where analysts must often infer the
behavior of code whose surface structure has been deliberately
engineered to resist comprehension.

Large language models (LLMs) offer a promising alternative for
automating such reasoning. Models have demonstrated strong
capabilities in code generation, code understanding, summarization, and vulnerability
detection~\cite{chen2021evaluatinglargelanguagemodels,li2022codebert,
guo2021graphcodebert,roziere2023codegen,
Competition-level-code-generation-with-AlphaCode}.
However, most evaluations are performed on naturally written programs,
where meaningful identifiers, familiar control structures, and common
coding idioms provide abundant lexical and structural cues. Obfuscated
programs remove or distort many of these cues while retaining the
underlying computation, creating a useful stress test of whether an
LLM actually reasons about program behavior or instead depends on
surface regularities.

Recent studies suggest that this distinction matters.
Rong {\em et al.}~\cite{ase26} showed that semantics-preserving
obfuscations can cause substantial degradation in LLM program
reasoning, including program-output prediction. Similarly, Nikiema
{\em et al.}~\cite{nikiema2025codebarrierllmsactually} reported a
statistically significant decline in model performance as the
complexity of obfuscation increases. Together, these findings indicate
that this limitation is especially problematic
for security applications: an AI-assisted malware analyzer should not
cease to understand a malicious computation merely due to obfuscation.

An important challenge, however, has received considerably less
attention: \emph{real obfuscated programs rarely contain only one
transformation}. Obfuscators can combine multiple techniques so that,
for example, renamed identifiers are embedded inside flattened control
flow, rewritten expressions, and additional irrelevant code. Such
{\bf \em stacked obfuscation} creates a fundamentally harder deployment
setting. Success on each transformation independently does not
necessarily imply success when those transformations occur together.
Consequently, simply demonstrating that a model can be adapted to
individual obfuscations is insufficient. For practical malware
analysis, the learned capability must also \emph{compose}.

% This leads to the central question studied in this paper:
% \emph{How should capabilities learned from individual obfuscations be
% composed so that an LLM can reason robustly about programs containing
% multiple transformations?} We investigate this question through four
% natural composition strategies operating at different levels. First,
% \emph{routing} trains a specialist for each transformation and attempts
% to select the appropriate specialist at inference time. Second,
% \emph{breadth training} exposes a single model to examples from many
% individual transformations, with the expectation that separately
% learned capabilities will recombine when transformations are stacked.
% Third, \emph{model merging} independently trains transformation
% specialists and then combines their learned parameter updates into a
% single model. Finally, \emph{paired-consistency anchoring} uses a
% clean-code teacher to regularize a student processing the corresponding
% obfuscated program through KL-divergence minimization. These strategies
% represent four different hypotheses about where composition should
% occur: in model selection, in the training distribution, in parameter
% space, or through behavioral alignment.

% Our empirical evaluation reveals substantial differences among these
% strategies. Routing provides little benefit because deciding which
% single specialist should process a stacked input does not effectively
% combine the knowledge required for multiple simultaneous
% transformations. Breadth training performs considerably better:
% training across diverse individual obfuscations enables the model to
% generalize to combinations of transformations represented by the
% training distribution. However, this advantage comes with an important
% cost. As breadth increases, the model becomes increasingly adapted to
% the transformation surfaces it has observed and can lose performance
% on a previously unseen obfuscation family. Thus, greater exposure to
% known transformations can improve reasoning over familiar stacks while
% simultaneously reducing robustness beyond them.

% Paired-consistency anchoring addresses this problem from a different
% direction. Rather than encouraging the model to learn increasingly many
% obfuscation surfaces, it constrains an obfuscated-code student to
% preserve the behavior of a model reasoning over the corresponding clean
% program. Specifically, for an obfuscated program $x_{\mathrm{obf}}$
% with clean parent $x_{\mathrm{clean}}$, the training objective combines
% the ordinary task loss with a KL-divergence term between the output
% distributions of the clean-code teacher and the obfuscated-code student:
% %
% \begin{equation}
% \mathcal{L}
% =
% \mathrm{CE}\!\left(y \mid x_{\mathrm{obf}}\right)
% +
% \lambda\,
% D_{\mathrm{KL}}\!\left(
% p_{\mathrm{teacher}}(\cdot \mid x_{\mathrm{clean}})
% \,\Vert\,
% p_{\mathrm{student}}(\cdot \mid x_{\mathrm{obf}})
% \right).
% \label{eq:paired_consistency}
% \end{equation}
% %
% This approach reduces the tendency of adaptation to drift toward
% obfuscation-specific behavior by explicitly tying the student's
% predictions to the semantics-equivalent clean program. Nevertheless,
% our evaluation shows that behavioral anchoring is not the strongest
% composition mechanism.

% Instead, \emph{model merging} provides the strongest overall
% performance. Unlike routing, which chooses among specialists, merging
% attempts to retain useful parameter updates from multiple specialists
% simultaneously. Unlike breadth training, it does not require a single
% optimization trajectory to learn all transformation surfaces jointly.
% And unlike paired-consistency anchoring, it is not constrained merely
% to reproduce the behavior of a clean-code teacher. By composing
% independently learned parameter updates, the merged model achieves
% stronger reasoning performance over stacked obfuscations than the
% routing, breadth-training, and KL-based alternatives.

%we propose composing independently trained obfuscation specialists using \emph{model merging}. We compare our approach against three baseline 

This leads to our central question:
\emph{How should capabilities learned from individual obfuscations be
composed so that an LLM can reason robustly about programs containing
multiple transformations without sacrificing capabilities that are not
reinforced during adaptation?} To address this question, we compare approaches that compose
obfuscation-specific knowledge at different stages of the pipeline:
breadth training, which exposes a single adapter to multiple transformations;
paired-consistency training, which regularizes an obfuscated-code student
against a clean-code teacher; inference-time routing, which maintains
separate specialists and combines them dynamically; and parameter-space
merging, which combines independently trained specialists before inference.

Our evaluation reveals a consistent trade-off between
\emph{specialization} and \emph{capability preservation}. Monolithic
fine-tuning methods achieve the highest peak accuracy when evaluation closely
matches the transformation distribution and reasoning direction used during
training. However, this specialization can come at a cost: performance
degrades more sharply when the model is evaluated on an unseen obfuscation
family or asked to reason in an untrained reverse direction. In particular,
fine-tuned models frequently lose capability that was present in the base
model even as their performance improves on the trained forward task.

In this paper, we propose composing independently trained obfuscation
specialists using {\bf model merging}. Rather than asking a single
optimization trajectory to internalize all transformations jointly, merging
first learns each transformation independently and then combines the resulting
parameter updates into one adapter. The resulting model requires neither an
expert bank nor an inference-time router. Although merged models do not always
achieve the highest in-distribution accuracy, they remain competitive on
familiar stacked obfuscations while preserving substantially more capability
outside the trained objective. Across the evaluated regimes, model merging
therefore achieves the strongest overall balance between performance on stacked
obfuscations and preservation of untrained reasoning capabilities.

This distinction becomes clear when the reasoning task is
reversed. Our adapters are trained only for forward program reasoning, yet we
also evaluate whether they can infer an input that produces a specified output.
Several monolithic fine-tuning approaches degrade below the untuned model on
this task, despite improving substantially on forward reasoning. Merged models,
by contrast, largely preserve or improve this pre-existing capability. We
further find that many fine-tuned failures are not arbitrary: the model often
returns the correct \emph{forward} answer to a backward prompt, indicating
strong specialization to the reasoning direction reinforced during training.

Merged models also substantially outperform their corresponding base models
on the original unobfuscated programs. Because the default merge includes a
clean-code specialist in addition to five obfuscation specialists, this result
by itself does not identify which specialists are responsible for the clean-code
gain. We therefore treat clean-code performance as evidence that merging can
preserve general program-reasoning capability while composing
obfuscation-specific adaptations, and separately examine specialist subsets to
test whether this capability also emerges from the obfuscation specialists
alone.

We hypothesize that the robustness of merging arises from the independence of
the specialist training trajectories. Semantics-preserving obfuscations provide
multiple \emph{views} of the same underlying computation. A specialist may
therefore learn both transformation-specific information associated with one
surface representation and updates that support reasoning about the computation
shared across representations. Training these specialists independently avoids
forcing all such signals through a single joint optimization trajectory.
Parameter-space composition can then combine complementary updates while
reducing the degree to which adaptation is dominated by one transformation
or reasoning direction. We test this hypothesis through clean-code evaluation,
specialist-subset ablations, unseen transformations, stacked obfuscations, and
reversed reasoning.

These results lead to a broader conclusion about robustness to code
obfuscation. The challenge is not simply to expose a model to more
transformations. Increasing training coverage can improve peak performance on
familiar conditions while simultaneously narrowing the set of capabilities
preserved by the adapted model. Robust adaptation therefore depends not only
on \emph{what} transformations are observed, but also on \emph{how} the
resulting knowledge is composed. Model merging provides a way to retain most
of the adaptation benefit while reducing this loss of untrained capability.


{\em Contributions.}
This paper makes the following main contributions:

\textbf{1. A systematic study of compositional robustness under stacked
obfuscation.}
We study whether capabilities learned from individual obfuscation
transformations compose when multiple transformations occur simultaneously,
reflecting an important characteristic of real obfuscations.

\textbf{2. An empirical characterization of specialization and forgetting
under obfuscation adaptation.}
We show that conventional fine-tuning achieves strong performance on familiar
transformations and stacked combinations, but can become increasingly
specialized to the observed transformation distribution and reasoning
direction. This specialization can degrade capabilities that are not reinforced
by the training objective, including reverse program reasoning.

\textbf{3. Model merging as a mechanism for balancing adaptation and
capability preservation.}
Parameter-space merging remains competitive with the strongest fine-tuning
approaches on stacked obfuscations while substantially better preserving
performance on an untrained reverse-reasoning task. Across evaluation regimes,
it provides the strongest overall balance between adaptation performance and
preservation of untrained reasoning capabilities, while producing a single
deployed adapter with no expert bank or inference-time router.

\textbf{4. Evidence that independently trained specialists preserve
complementary reasoning capabilities.}
Through clean-code evaluation, specialist-subset comparisons, unseen
transformations, and reversed reasoning, we examine whether independently
trained specialists retain capabilities shared across semantics-preserving
views in addition to transformation-specific adaptations.